\documentclass[12pt]{article}
\usepackage{enumitem}
\usepackage{sbc-template}
\usepackage{amsmath} 
\usepackage{graphicx,url,tabularx,booktabs}
\usepackage[english,brazil]{babel}
\usepackage[utf8]{inputenc}  
\usepackage{comment}
\usepackage{ifthen}
\usepackage{graphicx}
\usepackage{tcolorbox}
\usepackage{amssymb}
\usepackage{pifont}
\usepackage[table]{xcolor}
\usepackage{multirow} 
\newcommand{\cmark}{\textcolor{green!60!black}{\ding{51}}}
\newcommand{\xmark}{\textcolor{red!70!black}{\ding{55}}}
\usepackage{float}
\usepackage{breakcites}
\usepackage{makecell}
\usepackage{tikz}
\usetikzlibrary{arrows.meta,positioning,fit,backgrounds}
\sloppy

\title{AnonShield: Scalable On-Premise\\Pseudonymization for CSIRT Vulnerability Data}

\author{Cristhian Kapelinski\inst{1}, Douglas Lautert\inst{1}, Beatriz Machado\inst{1}\\Diego Kreutz\inst{1}, Isadora Garcia Ferrão\inst{2}}

\address{Federal University of Pampa (UNIPAMPA)\\\vspace{-8mm}
\nextinstitute
Université de Bretagne Occidentale (UBO)\\\vspace{-8mm}
 \email{\{cristhianavilla,douglaslautert,beatrizmachado\}.aluno@unipampa.edu.br}\\\vspace{-8mm}
   \email{diegokreutz@unipampa.edu.br, isadora/garciaferrao@univ-brest.fr}\\\vspace{-6mm}
}

\begin{document} 
\selectlanguage{english}
\maketitle
     
\begin{abstract}
We present AnonShield, a high-throughput, on-premise pseudonymization system that combines GPU-accelerated NER, streaming processing, caching, and schema-aware configuration. Evaluated on datasets up to 550\,MB (70{,}951 records), AnonShield reduces processing time from over 92 hours to under 10 minutes (up to 738$\times$ speedup) while achieving up to 94.2\% F1-score and 96.7\% recall.
Our results show that scalable pseudonymization of vulnerability data is feasible without sacrificing analytical utility, enabling compliant data sharing in operational CSIRT environments.
\end{abstract}

% \begin{resumo} 
% Relatórios de varredura de vulnerabilidades processados por CSIRTs frequentemente contêm Informações Pessoalmente Identificáveis (PII), criando desafios de conformidade com LGPD e GDPR. Este artigo apresenta o AnonShield, um framework de pseudonimização que combina NER acelerado por GPU, cache LRU, processamento em streaming e configuração orientada a esquema. Avaliado em datasets de até 550 MB (70.951 registros), reduz o tempo de processamento de ~92 horas para menos de 10 minutos (speedup de 738×). A melhor estratégia (filtered/hybrid) atinge F1 = 94,2\% e Recall = 96,7\% em um conjunto de validação anotado por especialistas.
% \end{resumo}

\section{Introduction}
\label{sec:intro}

Computer Security Incident Response Teams (CSIRTs) routinely process vulnerability scan reports generated by tools such as OpenVAS and Tenable across large and heterogeneous infrastructures, including containerized services, legacy systems, campus and backbone networks. These reports embed sensitive information, such as Personally Identifiable Information (PII), IP addresses, hostnames, and credentials, alongside technical findings~\cite{Kapelinski2025}. Sharing such data across organizations or machine learning pipelines exposes these identifiers to potential adversaries, raising compliance challenges under GDPR, LGPD, and internal data governance policies. Indeed, network identifiers like IP addresses may qualify as personal data under GDPR when linked to individuals~\cite{Albakri2019,Nweke2020}. As vulnerability datasets grow in volume, addressing this exposure at operational scale is critical for compliant data sharing and for enabling downstream uses such as LLM-based analysis and the generation of synthetic cyber threat data~\cite{almorjan2025synthetic}

The scale of this problem is rapidly increasing. In 2025, the National Vulnerability Database recorded 48{,}448 CVEs, a 20\% increase over 2024, and forecasts indicate approximately 59{,}400 CVEs in 2026~\cite{cvedetails2025,FIRST2026}. VulnCheck reported 884 Known Exploited Vulnerabilities in 2025, with nearly 29\% exploited at or before disclosure~\cite{VulnCheck2026}. In Brazil, CAIS/RNP continuously scans its infrastructure using Tenable, generating datasets exceeding 70{,}000 vulnerability records per cycle (Dataset~D2). Such volume makes manual analysis unfeasible and demands LLM-assisted processing, yet feeding raw operational data into such pipelines conflicts with GDPR/LGPD and data sovereignty constraints~\cite{ANPD2024}: LLM-based analysis requires context-rich but identifier-free data~\cite{Bandel2025}. Vulnerability scan data primarily raises security and sovereignty concerns; incident data additionally exposes PII of victims, adding a direct LGPD/GDPR compliance dimension. The tension between sharing vulnerability data and protecting embedded PII has been documented in Cyber Threat Intelligence literature~\cite{Albakri2019,Nweke2020,Wagner2016}, yet no widely adopted solution exists for large-scale vulnerability scan reports.

Table~\ref{tab:related_work} positions existing approaches across five dimensions relevant to CSIRT deployments. The solution space divides into three categories, each with structural limitations. Commercial cloud frameworks such as Google Cloud DLP~\cite{GoogleDLP} and Amazon Comprehend~\cite{AmazonComprehend} provide scalable pipelines with extensive entity coverage but require data transmission to external endpoints, violating data sovereignty constraints. Microsoft Presidio~\cite{MicrosoftPresidio}, while on-premise, lacks recognizers for cybersecurity-specific entities such as CVE identifiers, CPE strings, certificate serial numbers, and cryptographic artifacts, leaving parts of the vulnerability profile exposed. IRI DarkShield~\cite{DarkShield} exhibits similar limitations and introduces licensing restrictions that hinder adoption in academic and public-sector CSIRTs. Open-source tools address narrower problems. Anonip~\cite{Anonip} supports only IP anonymization, ARX~\cite{Prasser2014} focuses on tabular data and may degrade analytical utility through $k$-anonymity~\cite{Slijep2021}, and LogLicker~\cite{LogLicker2023} relies on static RegEx patterns that lead to higher false-negative rates in heterogeneous vulnerability data.

\begin{table}[!htp]
\centering
\caption{{\scriptsize Overview of anonymisation and pseudonymisation
tools and their primary gap relative to the CSIRT vulnerability
processing context. Data type: 
\textbf{US}~=~Unstructured; \textbf{S}~=~Structured.
\textbf{--} in Formats: tool operates on in-memory strings with no
native file I/O.}}
\label{tab:related_work}
\renewcommand{\arraystretch}{1.2}
\resizebox{\textwidth}{!}{%
\scriptsize
\begin{tabular}{|
    >{\raggedright\arraybackslash}p{2.2cm}|
    >{\raggedright\arraybackslash}p{3.4cm}|
    >{\raggedright\arraybackslash}p{1.9cm}|
    >{\raggedright\arraybackslash}p{0.7cm}|
    >{\raggedright\arraybackslash}p{2.8cm}|
    >{\raggedright\arraybackslash}p{4.2cm}|}
\hline
\textbf{System} &
\textbf{Main Technique} &
\textbf{Target Domain} &
\textbf{Type} &
\textbf{File Formats} &
\textbf{Gap in CSIRT context} \\
\hline

Microsoft Presidio~\cite{MicrosoftPresidio} &
NLP + RegEx + extensible recognizers &
Generic text & US &
-- &
No CVE/CPE/hash recognizers; domain-irrelevant patterns produce FPs
in vulnerability data \\
\hline

Google Cloud DLP~\cite{GoogleDLP} &
NLP + FPE + PII classification &
General data & US+S &
\texttt{.txt, .csv, .json,} \texttt{.pdf, .docx} &
Cloud endpoint conflicts with CSIRT data sovereignty policies;
no cybersecurity-specific entities \\
\hline

Amazon Comprehend~\cite{AmazonComprehend} &
Cloud NLP for PII detection and redaction &
Generic text & US &
\texttt{.txt} &
Cloud endpoint conflicts with CSIRT data sovereignty policies;
no vulnerability scanner format support \\
\hline

IRI DarkShield~\cite{DarkShield} &
Multi-source masking + REST API &
Corporate data & US+S &
\texttt{.txt, .pdf, .xml,} \texttt{.json, .sql, images} &
Commercial licence limits adoption in academic and public-sector
CSIRTs; no cybersecurity-specific entities \\
\hline

 
LogLicker~\cite{LogLicker2023} &
RegEx + masking + remapping manifest &
CloudTrail \& Generic logs & US+S &
\texttt{.json, .txt} &
RegEx only; lacks contextual NER; high manual maintenance for custom logs \\
\hline
 
Anonip~\cite{Anonip} &
Bit masking of trailing IP bits via pipe &
Web server logs & US &
\texttt{.log, .txt} &
Covers IP addresses only \\
\hline
 
ARX~\cite{Prasser2014} &
$k$-anonymity, $l, t, \delta$ privacy models &
Biomedical data & S &
\texttt{.csv, .xlsx, sql} &
Tabular data only; inapplicable to vulnerability reports \\
\hline
 
AnonLFI v1.0~\cite{Bandel2025} &
Hybrid NER + RegEx + SHA-256 (no salt) &
Incidents & US+S &
\texttt{.txt, .docx,} \texttt{.csv, .xlsx} &
Saltless SHA-256 vulnerable to rainbow-table attacks; no native XML/JSON support; impractical at scale \\
\hline

AnonLFI v2.0~\cite{Kapelinski2025} &
Hybrid NER + RegEx + OCR + HMAC-SHA256 + native XML/JSON processors &
Incidents \& Vulnerabilities & US+S &
\texttt{.xml, .json, .pdf,} \texttt{.txt, .csv, .docx,} \texttt{.xlsx,} images &
Fails at operational scale; superlinear latency from DOM parsing; no GPU or streaming; high variability (CV $>$ 0.96); $>$92 hours for 550\,MB \\
\hline

\textbf{AnonShield (this work)} &
\textbf{Hybrid NER + RegEx + OCR + HMAC-SHA256 + GPU + LRU cache + streaming} &
\textbf{Incidents \& Vulnerabilities} & \textbf{US+S} &
\texttt{.xml, .json, .pdf,} \texttt{.txt, .csv, .docx,} \texttt{.xlsx, .jsonl, .log,} images &
  \\
\hline
 
\end{tabular}}
\end{table}
 


Prior AnonLFI generations were the first solutions tailored to CSIRT contexts. AnonLFI v1.0~\cite{Bandel2025} introduced a hybrid NER and RegEx pipeline validated on 763 incidents, achieving Precision of 100\% and Recall of 97.38\%, but relies on saltless SHA-256 and lacks native support for XML and JSON. AnonLFI v2.0~\cite{Kapelinski2025} extended support to vulnerability data, adding HMAC-SHA256, native XML and JSON processing, and OCR capabilities, achieving F1-scores of 76.5\% for OCR and 92.13\% for OpenVAS XML. However, its estimated processing time exceeds 92 hours for operational datasets with 70{,}951 records, making it impractical at scale. Taken together, no existing solution simultaneously satisfies four key requirements: cryptographic pseudonymization with referential integrity, native processing of hierarchical formats, domain-specific entity recognition, and scalable on-premise execution.

This paper presents AnonShield, a  pseudonymization framework that addresses the limitations of prior frameworks and tools. AnonShield extends prior designs with GPU-accelerated NER, LRU caching, streaming I/O, and a schema-aware \texttt{anonymization\_config} mechanism that enables per-field policies without NER overhead. The contributions are threefold: (1) a scalable pseudonymization architecture for CSIRT vulnerability data, (2) a comparative evaluation of four strategies (\texttt{standalone}, \texttt{hybrid}, \texttt{presidio}, \texttt{filtered}) on datasets of up to 550\,MB (70{,}951 records), including comparisons with AnonLFI v1.0 and v2.0, and (3) an accuracy evaluation using a statistically representative sample of 67 vulnerability records annotated by three specialists. The evaluation spans both a controlled testbed with 130 services and an operational dataset from RNP infrastructure, demonstrating the feasibility of large-scale, on-premise pseudonymization.


%                       -
\section{The AnonShield Framework}
\label{sec:anonshield}
%                       -
AnonShield\footnote{\label{fn:anonshield_repo}\url{https://github.com/AnonShield/tool}} is a pseudonymisation framework for CSIRTs processing vulnerability scan reports and incident data. It detects and replaces sensitive entities with deterministic pseudonyms, preserving referential integrity across correlated reports. All processing is performed locally, ensuring compliance with GDPR~\cite{Amoo2024} and LGPD requirements.

The framework is organised into five sequential stages, illustrated in Figure~\ref{fig:pipeline}.

\begin{figure}[!htp]
    \centering
    \includegraphics[width=1\linewidth]{pipeline.png}
   \caption{AnonShield pipeline. The same entity produces a consistent pseudonym across documents given the same secret key.}
    \label{fig:pipeline}
\end{figure}

\textbf{(1)~Input.}
AnonShield supports XML, JSON, JSONL, CSV, TXT, LOG, PDF (text and image), DOCX, and XLSX formats. Large files are processed via incremental \textit{streaming}, using \texttt{ijson} for JSON and SAX parsing for XML, avoiding full in-memory loading.

\textbf{(2)~Entity Detection (NER).}
Detection combines a configurable transformer model, \texttt{Davlan/xlm-roberta-base-ner-hrl} (default, multilingual) or \texttt{attack-vector/SecureModernBERT-NER} (cybersecurity-focused), with 21 custom RegEx patterns covering IPv4/IPv6, URLs, hostnames, hashes, certificates, CVE/CPE identifiers, credentials, and others. Inference runs on GPU (CUDA), with an LRU cache to amortise repeated lookups across records.

\textbf{(3)~HMAC-SHA256 Pseudonymisation.}
Each entity is replaced by an HMAC-SHA256 slug computed with an operator-defined secret key. The slug length is configurable from 0 to 64 hexadecimal characters (up to 256\,bits of entropy), and the entity type is preserved as a prefix, e.g., \texttt{100.111.20.23}~$\to$~\texttt{[IP\_ADDRESS\_48624b5cdc]}.

\textbf{(4)~Persistence.}
The entity~$\leftrightarrow$~slug mapping is stored in a local SQLite database (\texttt{entity\_type}, \texttt{original\_name}, \texttt{slug\_name}, \texttt{full\_hash}, \texttt{first\_seen}, \texttt{last\_seen}), enabling controlled re-identification by an administrator holding both the database and the HMAC key. The database supports persistent and in-memory modes.

\textbf{(5)~Output.}
For structured formats (XML, JSON, JSONL, CSV, XLSX), the original schema and structure are preserved, with only sensitive values replaced. For PDF, DOCX, and image inputs, text is extracted via parsing or OCR, pseudonymised, and written to plain-text (\texttt{.txt}) output files.

AnonShield recognises 34 entity types across six categories, extending beyond classical PII to include cybersecurity-specific identifiers such as CVE IDs, CPE strings, UUIDs, and OIDs. Although publicly available, these identifiers may expose an organization’s vulnerability profile and enable targeted exploitation, making their sensitivity context-dependent. To address this, AnonShield decouples entity detection from replacement policy, enabling fine-grained control via \texttt{--entities-to-preserve} and \texttt{fields\_to\_exclude}, which support selective pseudonymization and field-level exclusion without altering the recognition pipeline.

AnonShield supports multiple anonymisation strategies, selectable via \texttt{-{}-anonymization-strategy}, to balance accuracy and performance. Presidio maximises recall but introduces false positives from domain-irrelevant recognizers. Filtered mitigates this by restricting recognizers to cybersecurity-relevant entities, while Hybrid preserves detection behaviour and reduces overhead through a lightweight replacement mechanism. Standalone bypasses Presidio entirely, achieving the highest throughput and lowest latency at a small recall cost due to span-boundary differences. An experimental SLM-based strategy, exploring locally deployed small language models for entity detection, is currently under development and remains outside the scope of this paper's evaluation. Additionally, the schema-aware \texttt{anonymization\_config} enables per-field control through \texttt{force\_anonymize}, \texttt{fields\_to\_anonymize}, and \texttt{fields\_to\_exclude}, improving both accuracy and performance and supporting reusable templates for formats such as STIX~2.x. Full details, including entity coverage, strategies, and configuration options, are available in the tool's GitHub repository\textsuperscript{\ref{fn:anonshield_repo}}.
 
 

\section{Experimental Methodology}
\label{sec:method}

The evaluation uses four datasets (Table~\ref{tab:datasets_summary}), covering
different scanning engines, operational scales, and processing challenges.

\begin{table}[!htp]
\centering
\caption{Dataset summary.}
\label{tab:datasets_summary}
\scriptsize
\begin{tabular}{lllrl}
\toprule
\textbf{ID} & \textbf{Source} & \textbf{Format} & \textbf{Size} & \textbf{Access} \\
\midrule
\textbf{D1}  & OpenVAS (VulnLab) & CSV/TXT/PDF/XML & 9.29/11.61/33.32/34.11\,MB & Public \\
\textbf{D1C} & Converted from D1 & XLSX/DOCX/JSON/PDF(img) & 3.34/7.58/60.72/1{,}479.56\,MB & Public \\
\textbf{D2}  & Tenable (CAIS/RNP)& CSV/JSON & 419.72/550.54\,MB & Private \\
\textbf{D3}  & Synthetic (Mock)  & CSV/JSON & 247.45/444.56\,MB & Public \\
\bottomrule
\end{tabular}
\end{table}

\textbf{D1 -- OpenVAS Heterogeneous Dataset.}
The VulnLab testbed\footnote{\url{github.com/CristhianKapelinski/LabVulnerabilities}}
comprises 130 targets across eight service categories (Table~\ref{tab:vulnlab_categories}),
producing 520 reports in four formats. The dataset also includes OpenVAS native
\texttt{.anonymous} XML files, which expose critical flaws in its built-in redaction.
While hostnames are masked, superficial sanitization leaves infrastructure artifacts
exposed, including IP addresses, environment names, and internal UUIDs. Most severely,
our empirical analysis found that 11 of the 130 reports explicitly leaked 71 unique
cryptographic hashes and TLS certificate fingerprints. Because these values are globally
indexed by search engines, they trivialize exact asset re-identification, rendering
the native anonymization fundamentally ineffective

\begin{table}[!htp]
\centering
\caption{VulnLab service categories.}
\label{tab:vulnlab_categories}
\scriptsize
\begin{tabular}{ll}
\hline
\textbf{Category} & \textbf{Examples} \\
\hline
Network \& Infrastructure  & DNS (Bind), FTP (ProFTPD), OpenLDAP \\
Web Applications \& APIs   & Juice Shop, DVWA, GraphQL \\
Web Servers \& Platforms   & Apache, Nginx, Tomcat, WordPress \\
Databases                  & MySQL, PostgreSQL, MongoDB, Redis \\
Monitoring \& Logging      & Elasticsearch, Prometheus, Grafana \\
DevOps \& CI/CD            & Jenkins, GitLab, Nexus \\
Messaging \& Streaming     & Kafka, RabbitMQ, Zookeeper \\
Operating Systems          & Ubuntu (14/16), Debian (8/9), CentOS (6/7) \\
\hline
\end{tabular}
\end{table}

\textbf{D1C -- Converted Formats.}
Each D1 report was converted to XLSX, DOCX, JSON, and image-only PDF, yielding
520 additional files. Rasterized PDFs inflate mean size by ${\approx}44\times$,
forcing the pipeline to rely exclusively on its OCR subsystem.

\textbf{D2 -- CAIS/RNP Operational Dataset.}
Restricted Tenable scans comprising 70{,}951 vulnerability records (JSON 550\,MB;
CSV 420\,MB), representative of large-scale enterprise environments and used as the
primary operational benchmark. Not publicly released due to privacy and security constraints. In fact, publishing vulnerability scan data is infeasible even after partial PII removal: retaining a vulnerability name or identifier (e.g., a CVE ID or plugin name) drastically reduces an attacker's search space. Removing all such identifiers to achieve safe publication would strip the dataset of its core content: a vulnerability dataset without vulnerability data has no analytical value. Consequently, for this data type, the most viable path to external sharing is synthetic data generation (an approach increasingly adopted for threat intelligence via LLMs \cite{almorjan2025synthetic}), while pseudonymization remains essential for enabling internal use of operational data in LLM pipelines without compromising organizational sovereignty or violating GDPR/LGPD. Pseudonymization and synthetic generation are thus complementary: anonymized data feeds LLMs that produce shareable synthetic datasets.

\textbf{D3 -- Synthetic Dataset.}
A public digital twin of D2, constructed by replacing each unique Tenable
\texttt{definition.id} with its corresponding CVE record from the official CVE
list,\footnote{\url{cve.org/downloads}} preserving D2's structural complexity and
entity volume while eliminating sensitive content.

\textbf{Hardware.}
All experiments ran on a single workstation: NVIDIA RTX 5060 Ti (16\,GB VRAM),
AMD Ryzen 5 8600G (6c/12t), 32\,GB DDR5-6000 RAM.

\textbf{Accuracy Evaluation.}
A statistically representative sample ($n{=}67$, population 6{,}472; 90\% CL,
$E{=}10\%$) was annotated by three domain specialists across 13 entity types,
yielding TP, FP, and FN counts for Precision, Recall, and F1-Score. AnonShield
used \texttt{SecureModernBERT-NER}; a preservation list (\texttt{TOOL},
\texttt{MALWARE}, \texttt{CVSS}) and \texttt{force\_anonymize} policies for
pattern-less fields (e.g., \texttt{Hostname}) were applied uniformly across all
versions under comparison.

\textbf{Statistical Analysis.}
A four-stage protocol was applied: Shapiro-Wilk normality test ($\alpha{=}0.05$);
Mann-Whitney U with Benjamini-Hochberg correction ($p_{\text{adj}}{<}0.05$);
Cohen's $d$ for effect size; and power-law regression ($T{=}a{\cdot}S^{\alpha}$)
combined with polynomial fitting to separate fixed overhead from throughput scaling.

\section{Results}
\label{sec:results}
% AnonShield is evaluated across three dimensions: (i) technical baseline and scalability (D1/D1C); (ii) operational throughput at scale (D2/D3); and (iii) pseudonymization accuracy. This structure validates the trade-off between PII mitigation and CSIRT analytical utility.

\subsection{Processing Performance}

\textbf{Heterogeneous formats (D1).}
Table~\ref{tab:d1_scientific_performance} (Appendix~\ref{sec:results_d1})
reports latency across 520 files ($N{=}260$ runs/strategy).
AnonShield\_standalone consistently achieves the lowest latency and highest
stability across all formats, with speedups of $16.51\times$ (XML),
$9.56\times$ (CSV), $3.27\times$ (PDF), and $3.04\times$ (TXT) over
AnonLFI~v2.0, all with large effect sizes (Cohen's $d \geq 0.42$).
AnonLFI~v2.0 exhibits pathological variability ($CV > 0.96$) caused by
memory-intensive DOM parsing; AnonShield reduces this to $CV < 0.8$ through
iterative streaming and GPU inference. Notably, AnonLFI~v2.0 regresses against
AnonLFI~v1.0 on CSV and TXT due to increased recognizer complexity without
architectural compensation, a limitation addressed by AnonShield.
Scalability analysis (Figure~\ref{fig:d1_scalability}, Appendix~\ref{sec:results_d1})
confirms near-linear complexity ($\alpha \approx 1$) for AnonShield, with an
amortization effect driving throughput above 60\,KB/s for larger files, while
AnonLFI~v2.0 degrades superlinearly.

\textbf{Converted formats (D1C).}
Results for XLSX, DOCX, JSON, and image-only PDF are detailed in
Table~\ref{tab:converted_perf} (Appendix~\ref{sec:results_d1c}).
AnonShield\_standalone achieves a $23.24\times$ speedup on JSON ($CV{=}0.60$),
$5.00\times$ on XLSX, and $2.08\times$ on DOCX. For image-only PDF, where the
${\approx}44\times$ size inflation shifts the bottleneck to OCR, speedup narrows
to $1.64\times$ while maintaining linear scaling ($R^2 \geq 0.989$). A single
file (\textit{openssh-server\_images.pdf}) caused 2 AnonShield failures due to
a malformed content stream triggering \texttt{PyMuPDF}'s strict parser, an
edge case absent in AnonLFI~v2.0 only because it lacks the corresponding
memory-management step.

\textbf{Operational scale (D2 and D3).}
Table~\ref{tab:large_scale_d2_merged} and Table~\ref{tab:large_scale_d3_merged}
(Appendix~\ref{sec:results_large_scale}) present the most consequential results
of this evaluation. On D2 (70{,}951 records, 550\,MB JSON), AnonShield\_standalone
completes processing in $453\,\text{s}$ ($1{,}250\,\text{KB/s}$), whereas
AnonLFI~v2.0's estimated runtime exceeds ${\sim}$92 hours, a speedup of
${\geq}738\times$. Activating the schema-aware \texttt{anonymization\_config}
reduces D2-CSV runtime further to $12.55\,\text{s}$ ($34{,}341\,\text{KB/s}$),
a $46.9\times$ gain over full NER inference, by enabling deterministic
field-level remapping that bypasses the inference pipeline entirely.
On D3, GPU acceleration yields speedups of up to $3{,}532\times$ over
AnonLFI~v2.0; even on CPU-only hardware, AnonShield\_standalone achieves
${\geq}535\times$ speedup, confirming that performance gains are architectural
rather than hardware-dependent. Throughput on D3 ($3{,}473\,\text{KB/s}$, CSV)
exceeds D2 ($732\,\text{KB/s}$) due to higher entity redundancy enabling
effective LRU cache reuse; D2's lower cache hit rate, driven by high-entropy
operational data, reflects realistic worst-case deployment conditions.
All strategy comparisons are statistically significant ($p_{\text{adj}} < 0.001$,
Mann-Whitney U with Benjamini-Hochberg correction).

\subsection{Accuracy}

Table~\ref{tab:accuracy} presents results on the 67-record specialist-annotated
validation set (13 entity types). Recall is the operationally critical metric:
an unredacted False Negative exposes sensitive infrastructure details to adversaries and constitutes a potential GDPR/LGPD breach, whereas a False Positive degrades analytical utility without compromising security or compliance
Accuracy improves substantially across generations: AnonLFI~v1.0 achieves
F1~$=$~23.7\% (Recall~$=$~18.8\%); AnonLFI~v2.0 reaches 54.0\%
(Recall~$=$~40.7\%); and AnonShield \texttt{filtered}/\texttt{hybrid} achieve
F1~$=$~94.2\% at Recall~$=$~96.7\%. The \texttt{presidio} strategy matches this
Recall but at lower Precision (71.6\%), favoring high-sensitivity environments
where over-anonymization is acceptable. \texttt{Standalone} trails by 2.2\,pp
in Recall (94.5\%) at the highest throughput, a gap attributable to
span-boundary behavior rather than a fundamental model limitation.

All 25 common False Negatives across AnonShield strategies originate from two
sources: (1)~19 unrecognized organization names
(\textit{Oracle}, \textit{ISC}, \textit{Wiesemann \& Theis}) in formulaic
vulnerability attribution strings, where entity context is insufficient for
NER disambiguation; and (2)~6 structural artifacts in semi-structured
free-text, including port values appearing as plain integers
(e.g., \texttt{80 or 443}), where neither regex nor contextual NER can infer
entity type without schema information. These failure modes are deterministic
and enumerable, making them addressable via targeted recognizer extensions.
False Positives are dominated by Presidio's \texttt{DATETIME} and
\texttt{IP\_ADDRESS} recognizers misclassifying version strings
(e.g., \texttt{2.4.51}~$\to$~\texttt{[DATE\_TIME\_...]}); \texttt{filtered}
and \texttt{hybrid} suppress this to 64 FPs via a curated recognizer subset.
A systematic AnonLFI~v2.0 misclassification of \texttt{Severity:~Medium} as
\texttt{ORGANIZATION} is fully resolved in AnonShield via
\texttt{anonymization\_config}. Full False Negative and False Positive
breakdowns are provided in Appendix~\ref{sec:accuracy_results}.



\section{Availability, Implementation and Demonstration Plan}
\label{sec:availability}

AnonShield is publicly available on GitHub\textsuperscript{\ref{fn:anonshield_repo}}, including source code, documentation, datasets, and all evaluation artefacts. The repository provides step-by-step installation instructions, a minimal functional test executable in under 5 minutes, and scripts to reproduce the main experimental results. The tool is implemented in Python (3.12), uses Microsoft Presidio for NER orchestration, and is managed via \texttt{uv}, with no external service dependencies. It is also distributed as a Docker image in CPU and GPU (CUDA) variants on Docker Hub\footnote{\url{https://hub.docker.com/r/anonshield/anon}}.

For the SBRC Tools Lounge, the demonstration will run on a standard laptop (4 vCPUs, 8 GB RAM) running Ubuntu, without requiring specialized hardware or network infrastructure. A GPU is optional, as the CPU version is fully functional. The demonstration includes live pseudonymization of OpenVAS reports, comparison of anonymization strategies, and illustration of throughput gains enabled by the schema-aware \texttt{anonymization\_config} mechanism.

\section{Conclusion and Future Work}
\label{sec:conclusion}

This paper presented AnonShield, the third generation of the AnonLFI research line, achieving up to ${\sim}738\times$ speedup (from ${\sim}$92.9 hours to under 10 minutes on 550\,MB) and ${\geq}535\times$ on CPU-only hardware. The \texttt{filtered} and \texttt{hybrid} strategies deliver the best accuracy-performance trade-off (F1~=~94.2\%, Recall~=~96.7\%), while \texttt{anonymization\_config} provides additional gains of up to $47\times$. These results demonstrate that large-scale, on-premise pseudonymization is operationally practical for CSIRT workflows.

Two boundaries define the current scope. PDF and image accuracy is constrained by text extraction: PyMuPDF and OCR introduce spacing and formatting breaks that fragment entities, and malformed PDFs trigger hard failures rather than graceful fallbacks. We are also working on extracting vulnerabilities from scanner reports as structured datasets for direct use in AnonShield~\cite{machado2025}. Formal privacy models (k-anonymity, l-diversity, t-closeness, differential privacy~\cite{sweeney2002k, machanavajjhala2007l, li2007t}) are most applicable to incident data, where victim PII introduces quasi-identifier risks beyond the security and sovereignty concerns common to both data types; in vulnerability contexts, publishing raw records is infeasible regardless of anonymization depth, as discussed above.

Future work targets improved PDF robustness, refinement of \texttt{standalone}, completion of the SLM-based strategy, extended evaluations, and modelling privacy-utility trade-offs. Together, these advances aim to establish privacy-preserving, analytically rich data exchange as standard practice across CSIRT networks at a time when the global vulnerability landscape is accelerating faster than human analysis can track.

\bibliographystyle{sbc}
\begingroup
\makeatletter
\let\oldthebibliography\thebibliography
\let\endoldthebibliography\endthebibliography
\renewcommand{\thebibliography}[1]{%
  \oldthebibliography{#1}%
  \setlength{\itemsep}{-0.5ex}%
  \setlength{\parskip}{4pt}%
  \setlength{\parsep}{0pt}%
  \setlength{\topsep}{0pt}%
}
\makeatother
\bibliography{sbc-template}
\endgroup

% ============================================================
% APPENDIX: RESEARCH ARTEFACTS
% ============================================================
\appendix

\section{Performance and Scaling Analysis (Dataset D1)}
\label{sec:results_d1}

Performance metrics for 520 files ($N=260$ runs/strategy) are summarized in Table~\ref{tab:d1_scientific_performance}. Due to file-size heterogeneity and right-skewed distributions (Shapiro-Wilk $p < 10^{-6}$), the \textbf{median} is used as a robust measure for D1, while the \textbf{mean} ($\bar{t} \pm \sigma$) is reserved for D2/D3.

%   TABELA (Preservada conforme solicitado)  
\begin{table}[!htp]
\centering
\caption{Processing latency on D1 reports.}
\label{tab:d1_scientific_performance}
\scriptsize
\begin{tabular}{llrrrrrrc}
\toprule
\textbf{Format} & \textbf{Version/Strategy} & \textbf{Mean (s)} & \textbf{Median (s)} & \textbf{Max (s)} & \textbf{CV} & \textbf{Speedup} & \textbf{Cohen's $d$} \\
\midrule
\multirow{3}{*}{\textbf{XML}} & AnonLFI~v2.0 (Baseline) & 191.89 & 176.44 & 1717.71 & 0.96 & $1.00\times$ & -- \\
 & AnonShield\_presidio & 15.62 & 13.69 & 85.58 & 0.60 & $12.28\times$ & 1.29 (Large) \\
 & \textbf{AnonShield\_standalone} & \textbf{11.62} & \textbf{10.02} & \textbf{72.34} & \textbf{0.63} & $\mathbf{16.51\times}$ & \textbf{1.32 (Large)} \\
\midrule
\multirow{4}{*}{\textbf{CSV}} & AnonLFI~v2.0 (Baseline) & 74.18 & 37.01 & 1442.08 & 1.75 & $1.00\times$ & -- \\
 & AnonLFI~v1.0 & 37.92 & 20.41 & 577.96 & 1.55 & $1.96\times$ & 0.35 (Small) \\
 & AnonShield\_presidio & 10.65 & 9.08 & 60.62 & 0.53 & $6.97\times$ & 0.67 (Medium) \\
 & \textbf{AnonShield\_standalone} & \textbf{7.76} & \textbf{6.67} & \textbf{47.87} & \textbf{0.53} & $\mathbf{9.56\times}$ & \textbf{0.70 (Medium)} \\
\midrule
\multirow{3}{*}{\textbf{PDF}} & AnonLFI~v2.0$^{\dag}$ (Baseline) & 27.38 & 13.01 & 304.24 & 1.70 & $1.00\times$ & -- \\
 & AnonShield\_presidio & 11.20 & 9.25 & 68.13 & 0.58 & $2.44\times$ & 0.47 (Small) \\
 & \textbf{AnonShield\_standalone} & \textbf{8.36} & \textbf{6.91} & \textbf{50.78} & \textbf{0.60} & $\mathbf{3.27\times}$ & \textbf{0.56 (Medium)} \\
\midrule
\multirow{4}{*}{\textbf{TXT}} & AnonLFI~v2.0 (Baseline) & 31.23 & 12.61 & 747.74 & 2.20 & $1.00\times$ & -- \\
 & AnonLFI~v1.0 & 20.28 & 10.98 & 354.91 & 1.56 & $1.54\times$ & 0.19 (Negligible) \\
 & AnonShield\_presidio & 13.10 & 10.21 & 100.28 & 0.73 & $2.38\times$ & 0.36 (Small) \\
 & \textbf{AnonShield\_standalone} & \textbf{10.28} & \textbf{7.91} & \textbf{80.71} & \textbf{0.78} & $\mathbf{3.04\times}$ & \textbf{0.42 (Small)} \\
\bottomrule
%\multicolumn{4}{l}
\end{tabular}
{\scriptsize $^{\dag}$~AnonLFI~v2.0 PDF failed runs excluded due to out-of-memory errors. $^{\dag}$$^{\dag}$~Presidio-based strategies are statistically similar ($p_{adj} > 0.46$), with AnonShield\_presidio as representative. AnonLFI~v1.0 does not support XML or PDF.}
\end{table}

The \texttt{standalone} strategy consistently outperforms Presidio-based versions ($1.3\times$--$1.4\times$ faster) by bypassing orchestration overhead.

Further scalability insights are illustrated in Figure~\ref{fig:d1_scalability}. Panel C highlights an \textbf{amortization effect}: as file size increases, AnonShield\_standalone throughput rises (peaking $>60$~KB/s), effectively spreading fixed startup costs. Conversely, AnonLFI~v2.0 shows degrading throughput, confirming its inefficiency for operational-scale workloads.

\begin{figure}[!htp]
\centering
\includegraphics[width=0.9\textwidth]{images/14_scaling_simplified.pdf}
\caption{Scalability analysis on D1 PDF subset: (A) Linear scale, (B) Log-Log complexity ($\alpha \approx 1$ for AnonShield), and (C) Throughput stability.}
\label{fig:d1_scalability}
\end{figure}

\section{Converted Format Benchmarks (D1C)}
\label{sec:results_d1c}

To evaluate versatility, 130 reports were converted into XLSX, DOCX, JSON, and image-only PDF ($N=260$ runs/strategy). Table~\ref{tab:converted_perf} summarizes the results, with Kruskal-Wallis tests confirming significant strategy differences across all formats ($p < 10^{-5}$).

\textbf{JSON and Office Formats.} In XLSX and DOCX, AnonLFI~v2.0 showed a \textbf{performance regression} (slower than AnonLFI~v1.0) because its additional recognizers lacked the architectural optimizations of AnonShield.

\textbf{Image-only PDF and Specific Failures.} The 2 AnonShield failed runs stem from \texttt{PyMuPDF}'s \texttt{clean\_contents()} function for memory management: this stricter parser triggered a \texttt{RuntimeError} (\textit{FzErrorArgument}) upon encountering a malformed content stream in the source PDF. AnonLFI~v2.0 bypassed this error as it lacked this cleaning step.


\begin{table}[!htp]
\centering
\caption{Processing latency on D1-converted formats.}
\label{tab:converted_perf}
\scriptsize
\begin{tabular}{llrrrrrr}
\toprule
\textbf{Format} & \textbf{Version/Strategy} & \textbf{Mean (s)} & \textbf{Max (s)} & \textbf{CV} & \textbf{Speedup} & \textbf{Cohen's $d$} \\
\midrule
\multirow{3}{*}{\textbf{JSON}}
 & AnonLFI~v2.0 (Baseline) & 246.55 & 1861.01 & 1.16 & $1.00\times$ & -- \\
 & AnonShield\_presidio            &  20.15 &   110.78 & 0.72 & $12.23\times$ & 1.08 (Large) \\
 & \textbf{AnonShield\_standalone}&  \textbf{10.61} & \textbf{49.02} & \textbf{0.60} & $\mathbf{23.24\times}$ & \textbf{1.12 (Large)} \\
\midrule
\multirow{4}{*}{\textbf{XLSX}}
 & AnonLFI~v1.0 (Baseline) & 35.51 & 320.73 & 1.37 & $1.00\times$ & -- \\
 & AnonLFI~v2.0            & 60.30 & 596.95 & 1.52 & $0.59\times$ & $-0.34$ (Small) \\
 & AnonShield\_presidio            &  9.86 &  38.64 & 0.47 & $3.60\times$ & 0.73 (Medium) \\
 & \textbf{AnonShield\_standalone}&  \textbf{7.11} & \textbf{27.95} & \textbf{0.47} & $\mathbf{5.00\times}$ & \textbf{0.82 (Large)} \\
\midrule
\multirow{4}{*}{\textbf{DOCX}}
 & AnonLFI~v1.0 (Baseline) & 19.58 & 186.55 & 1.34 & $1.00\times$ & -- \\
 & AnonLFI~v2.0            & 29.90 & 431.60 & 1.94 & $0.65\times$ & $-0.23$ (Small) \\
 & AnonShield\_presidio            & 12.16 &  63.24 & 0.68 & $1.61\times$ & 0.36 (Small) \\
 & \textbf{AnonShield\_standalone}&  \textbf{9.43} & \textbf{51.51} & \textbf{0.73} & $\mathbf{2.08\times}$ & \textbf{0.52 (Medium)} \\
\midrule
\multirow{3}{*}{\textbf{PDF-image$^{\dag}$}}
 & AnonLFI~v2.0 (Baseline) & 59.15 & 757.63 & 1.85 & $1.00\times$ & -- \\
 & AnonShield\_presidio            & 37.85 & 362.81 & 1.46 & $1.56\times$ & 0.23 (Small) \\
 & \textbf{AnonShield\_standalone}& \textbf{36.04} & \textbf{348.63} & \textbf{1.54} & $\mathbf{1.64\times}$ & \textbf{0.27 (Small)} \\
\bottomrule
\end{tabular}

{\scriptsize
$^{\dag}$ Failures: 1 file (\textit{openssh-server\_images.pdf}) failed in both runs ($N=2$) across all AnonShield strategies. $^{\dag}$$^{\dag}$\texttt{AnonShield\_presidio} represents the Presidio cluster. Speedup for JSON/PDF-image vs. AnonLFI~v2.0; XLSX/DOCX vs. AnonLFI~v1.0.}
\end{table}

\section{Large-Scale Operational Benchmarks (D2 and D3)}
\label{sec:results_large_scale}

AnonShield was evaluated against operational datasets D2 (Tenable, 550\,MB) and D3 (Mock CVE, 444\,MB). Due to the prohibitive slowness of AnonLFI v1.0 and v2.0, their runtimes were estimated based on D1 throughput. For CSV, these are \textbf{conservative lower bounds} (${\geq}121$ hours), as v2.0’s superlinear scaling significantly inflates costs as file sizes grow. For D3, an additional CPU-only benchmark (no GPU) was conducted to assess whether AnonShield's throughput gains depend on GPU hardware.

%   TABELA D2 (Completa)  
\begin{table}[!htp]
\centering
\caption{D2 (Operational) processing times (70,951 records).}
\label{tab:large_scale_d2_merged}
\scriptsize
\begin{tabular}{llrrrr}
\toprule
 & & \multicolumn{2}{c}{\textbf{CSV (419.72\,MB)}} & \multicolumn{2}{c}{\textbf{JSON (550.54\,MB)}} \\
\cmidrule(lr){3-4}\cmidrule(lr){5-6}
\textbf{Mode} & \textbf{Version / Strategy} & \textbf{Time (s)} & \textbf{KB/s} & \textbf{Time (s)} & \textbf{KB/s} \\
\midrule
\multirow{3}{*}{\textbf{No config}}
 & AnonLFI~v2.0 \emph{(est.)} & ${\geq}437{,}335$ & 0.98 & ${\sim}334{,}472$ & 1.69 \\
 & AnonShield\_presidio & $2{,}520.7 \pm 68.0$ & 171 & $985.7 \pm 73.7$ & 575 \\
 & \textbf{AnonShield\_standalone} & $\mathbf{588.5 \pm 30.7}$ & \textbf{732} & $\mathbf{453.1 \pm 35.9}$ & \textbf{1{,}250} \\
\midrule
\multirow{2}{*}{\textbf{With config}}
 & AnonShield\_presidio & $13.42 \pm 0.13$ & 32{,}034 & $18.88 \pm 0.16$ & 29{,}855 \\
 & \textbf{AnonShield\_standalone} & $\mathbf{12.55 \pm 0.74}$ & \textbf{34{,}341} & $\mathbf{18.03 \pm 0.23}$ & \textbf{31{,}272} \\
\midrule
\multicolumn{2}{l}{\textbf{Speedup standalone vs. v2.0 (no config)}} & \multicolumn{2}{c}{$\mathbf{{\geq}743\times}$} & \multicolumn{2}{c}{$\mathbf{{\sim}738\times}$} \\
\bottomrule
\end{tabular}

\scriptsize AnonShield\_standalone is significantly faster than Presidio-based strategies ($p_{adj} < 0.001$).
\end{table}

\textbf{Analysis of Large-Scale Impact.}
Without schema-aware configuration (full NER inference active), the GPU variant is up to ${\sim}6.6\times$ faster than CPU on CSV and ${\sim}5.1\times$ on JSON; with config (inference bypass), both perform similarly (${\sim}8\text{--}9$\,s).

%   TABELA D3 (Completa com Presidio)  
\begin{table}[!htp]
\centering
\caption{D3 (Synthetic) processing times.}
\label{tab:large_scale_d3_merged}
\scriptsize
\begin{tabular}{llrrrr}
\toprule
 & & \multicolumn{2}{c}{\textbf{CSV (247.45\,MB)}} & \multicolumn{2}{c}{\textbf{JSON (444.56\,MB)}} \\
\cmidrule(lr){3-4}\cmidrule(lr){5-6}
\textbf{Mode} & \textbf{Version / Strategy} & \textbf{Time (s)} & \textbf{KB/s} & \textbf{Time (s)} & \textbf{KB/s} \\
\midrule
\multirow{3}{*}{\textbf{No config (GPU)}}
 & AnonLFI~v2.0 \emph{(est.)} & ${\geq}257{,}835$ & 0.98 & ${\sim}270{,}086$ & 1.69 \\
 & AnonShield\_presidio & $318.0 \pm 5.1$ & 797 & $339.6 \pm 14.4$ & 1{,}343 \\
 & \textbf{AnonShield\_standalone} & $\mathbf{73.0 \pm 1.6}$ & \textbf{3{,}473} & $\mathbf{172.1 \pm 6.2}$ & \textbf{2{,}647} \\
\midrule
\textbf{No config (CPU)} & \textbf{AnonShield\_standalone} & $\mathbf{481.5 \pm 8.9}$ & \textbf{526} & $\mathbf{881.9 \pm 57.7}$ & \textbf{518} \\
\midrule
\multirow{2}{*}{\textbf{With config}}
 & AnonShield\_presidio & $8.89 \pm 0.06$ & 28{,}517 & $21.11 \pm 0.22$ & 21{,}571 \\
 & \textbf{AnonShield\_standalone} & $\mathbf{7.96 \pm 0.08}$ & \textbf{31{,}856} & $\mathbf{20.43 \pm 0.81}$ & \textbf{22{,}314} \\
\midrule
\multicolumn{2}{l}{\textbf{Speedup standalone (GPU) vs. v2.0}} & \multicolumn{2}{c}{$\mathbf{{\geq}3{,}532\times}$} & \multicolumn{2}{c}{$\mathbf{{\sim}1{,}569\times}$} \\
\multicolumn{2}{l}{\textbf{Speedup standalone (CPU) vs. v2.0}} & \multicolumn{2}{c}{$\mathbf{{\geq}535\times}$} & \multicolumn{2}{c}{$\mathbf{{\sim}306\times}$} \\
\bottomrule
\end{tabular}

{\scriptsize Standalone vs. Presidio comparisons are statistically significant in all cases ($p_{adj}<0.001$).}
\end{table}


\section{Accuracy Evaluation}
\label{sec:accuracy_results}

Partial anonymizations were counted as 1\,TP + 1\,FN. For \texttt{standalone}, all counts were computed against actual occurrences in the original text rather than the pseudonym count in the output: lacking an overlap-resolution layer, \texttt{standalone} can generate more pseudonyms than real occurrences (e.g., different recognizers detecting overlapping spans independently), which would artificially inflate TP and FP counts if the output were used as the counting basis. The \texttt{standalone} Recall gap (2.2\,pp) is a mechanical artifact of this same span-boundary behavior: its sequential replacement occasionally truncates entities (e.g., \texttt{[IP\_ADDRESS\_...].1.72}) or merges overlapping patterns.

\begin{table}[!htp]
\centering
\caption{Accuracy on 67-record validation set (D1 OpenVAS CSV).}
\label{tab:accuracy}
\small
\resizebox{0.6\linewidth}{!}{
\begin{tabular}{lrrrrrrr}
\toprule
\textbf{Version/Strategy} & \textbf{TP} & \textbf{FP} & \textbf{FN}
  & \textbf{Prec.} & \textbf{Rec.} & \textbf{F1} \\
\midrule
AnonLFI v1.0               & 106 & 224 & 459 & 32.1\% & 18.8\% & 23.7\% \\
AnonLFI v2.0               & 279 &  68 & 407 & 80.4\% & 40.7\% & 54.0\% \\
\midrule
AnonShield\_presidio   & 724 & 287 & 25  & 71.6\% & 96.7\% & 82.3\% \\
AnonShield\_filtered   & 724 &  64 & 25  & 91.9\% & 96.7\% & 94.2\% \\
AnonShield\_hybrid     & 724 &  64 & 25  & 91.9\% & 96.7\% & 94.2\% \\
AnonShield\_standalone & 739 & 102 & 43  & 87.9\% & 94.5\% & 91.1\% \\
\bottomrule
\end{tabular}
}
\end{table}

\textbf{False Negatives.} All AnonShield strategies share 25 common FNs from two
main sources: (1)~19~FN from organization names \emph{Oracle}~(13),
\emph{ISC}~(5), and \emph{Wiesemann \& Theis}~(1) unrecognized by
SecureModernBERT-NER in formulaic vulnerability attribution strings; and
(2)~6~FN from structural artifacts embedded in semi-structured free-text
fragments undetectable by both regex and NER, including partially anonymized
file paths, a URL, a hostname, and port values appearing as plain integers
in free-text (e.g., \texttt{80 or 443})---a case where neither pattern
matching nor contextual NER can infer the entity type without schema
information. Notably, the \texttt{xlm-roberta} model used in AnonLFI v1.0/v2.0
detected \emph{Oracle} inconsistently: the same entity could be anonymized
in one field and leaked in another within the same record, illustrating
complementary strengths across transformer architectures. \textbf{Anonymizing
$k{-}1$ out of $k$ occurrences is functionally equivalent to anonymizing
none}, since a single leaked instance suffices for cross-reference
re-identification.

\textbf{False Positives.} \texttt{AnonShield\_presidio} produces 287 FPs,
primarily from \texttt{DATETIME}/\texttt{IP\_ADDRESS} recognizers
misclassifying version strings (e.g., \texttt{2.4.51}~$\to$~\texttt{[DATE\_TIME\_...]})
and domain-irrelevant patterns (e.g., \texttt{UK\_NHS} triggering on RFC timestamps).
\texttt{Filtered} and \texttt{hybrid} reduce this to 64 via a curated
recognizer subset that suppresses low-relevance patterns.
\texttt{Standalone}'s 102~FPs stem additionally from entity-merging artifacts
and CVE regex splits that break referential integrity
(e.g., \texttt{CVE-2014-3597}~$\to$~\texttt{-[CVE\_ID\_...]VE-2014-3597}),
both absent in Presidio-based strategies due to the \texttt{AnonymizerEngine}
conflict-resolution pass. A systematic FP from AnonLFI~v2.0---the NER model classifying
\texttt{Severity: Medium} as an \texttt{ORGANIZATION}---is fully resolved in AnonShield via \texttt{anonymization\_config}.


% \section{Initialization Overhead}
% \label{sec:results_overhead}

% Table~\ref{tab:overhead} isolates the fixed startup cost model loading and
% pipeline instantiation by running all strategies on a near-empty 5-byte file
% (10 runs each). This overhead is paid once per invocation and is independent
% of file size.

% \begin{table}[!htp]
% \centering
% \caption{Initialization overhead measured on a 5-byte file (10 runs,
% $\bar{t} \pm \sigma$, seconds).}
% \label{tab:overhead}
% \small
% \begin{tabular}{lr}
% \toprule
% \textbf{Version/Strategy} & \textbf{Mean Time (s)} \\
% \midrule
% AnonLFI~v1.0             & $6.29 \pm 0.05$ \\
% AnonLFI~v2.0             & $7.22 \pm 0.06$ \\
% AnonShield\_presidio            & $6.83 \pm 0.06$ \\
% AnonShield\_hybrid              & $6.83 \pm 0.06$ \\
% AnonShield\_filtered     & $6.83 \pm 0.04$ \\
% \textbf{AnonShield\_standalone} & $\mathbf{4.93 \pm 0.03}$ \\
% \bottomrule
% \end{tabular}
% \end{table}

% The \texttt{standalone} strategy has the lowest startup cost at 4.93\,s,
% which is 32\% lower than AnonLFI~v2.0.
% The three Presidio-based AnonShield strategies share the same overhead (6.83\,s)
% because they all instantiate the Presidio AnalyzerEngine at startup.
% \texttt{Standalone} avoids this by loading only the models it needs directly,
% without the Presidio orchestration layer.


% \section{Limitations}
% \label{sec:disc_limitations}

% Several limitations should be noted. Accuracy was evaluated on OpenVAS CSV data from a controlled testbed, and the 67-record sample limits estimates for rare entity types; operational datasets may exhibit greater variability. Performance results for AnonLFI v1.0/v2.0 on large datasets (D2/D3) are throughput-based extrapolations, thus reported speedups represent conservative lower bounds. The SLM-based strategy was not evaluated, and \texttt{force\_anonymize} does not prevent leakage in free-text fields lacking identifiable patterns. Additionally, Precision, Recall, and F1 are dataset- and policy-dependent and should be interpreted as relative indicators rather than absolute guarantees. The Recall gap of \texttt{standalone} compared to Presidio-based strategies is due to localized implementation issues, not architectural limitations, and is expected to be addressed in future work.

\section{Research Artefacts}
\label{sec:artefacts}
Dataset D2 is restricted by a CAIS/RNP agreement. All other artefacts, reproducibility scripts, and claims verification are publicly available at:
\url{https://github.com/AnonShield/tool/blob/main/paper_data/AVAILABILITY.md}.

\end{document}